% =====================================================================
%  Báo cáo kỹ thuật hệ thống Marketplace Price Intelligence
%  Biên dịch: xelatex bao-cao-ky-thuat.tex (2 lần)  hoặc  tectonic bao-cao-ky-thuat.tex
%  Hình: đặt file vào figures/<tên>.pdf|png|jpg, tên file in ở góc mỗi ô trống.
% =====================================================================
\documentclass[11pt,a4paper]{article}
\newcommand{\ReportTitle}{Marketplace Price Intelligence}
\newcommand{\ReportPeriod}{Báo cáo kỹ thuật}
\newcommand{\StudentName}{Dương Đức Duy}
\newcommand{\StudentId}{20225828}
\newcommand{\Supervisor}{Trần Văn Đặng}
\newcommand{\Baseline}{nhánh \texttt{docs/system-report}, commit \texttt{93d7ea3}}
% =====================================================================
%  weekly-preamble.tex — khung cho báo cáo kỹ thuật
%  Biên dịch bằng XeLaTeX (hoặc tectonic). Không dùng pdfLaTeX.
% =====================================================================
% Lề theo chuẩn văn bản: trên/dưới 25 mm, trái 30 mm (đóng gáy), phải 20 mm
\usepackage[a4paper,top=25mm,bottom=25mm,left=30mm,right=20mm,
  headheight=14pt,footskip=12mm]{geometry}
\usepackage{fontspec}
\setmainfont{texgyrepagella}[Extension=.otf, UprightFont=*-regular, BoldFont=*-bold,
  ItalicFont=*-italic, BoldItalicFont=*-bolditalic]
\setsansfont{texgyreheros}[Extension=.otf, UprightFont=*-regular, BoldFont=*-bold,
  ItalicFont=*-italic, BoldItalicFont=*-bolditalic]
\setmonofont{lmmono10}[Extension=.otf, UprightFont=*-regular, ItalicFont=*-italic,
  BoldFont=lmmonolt10-bold, Scale=0.92]
\usepackage{polyglossia}
\setdefaultlanguage{vietnamese}

\usepackage{amsmath,amssymb}
\usepackage{xcolor}
\usepackage{booktabs,array,xltabular}
\usepackage{etoolbox,needspace}
\usepackage[section]{placeins}  % hình không trôi sang mục sau
\usepackage{enumitem}
\usepackage{fancyvrb}
\usepackage{titlesec}
\usepackage{fancyhdr}
\usepackage{tcolorbox}
\tcbuselibrary{skins,breakable}
\usepackage[hidelinks,unicode]{hyperref}

\setlength{\parindent}{0pt}
\setlength{\parskip}{4pt plus 1pt}
\setlength{\emergencystretch}{3em}
\tolerance=2000
\renewcommand{\arraystretch}{1.2}
\setlength{\tabcolsep}{5pt}
% Bảng: chữ nhỏ, sát lề trái, khoảng cách trên/dưới đều nhau
\setlength{\LTleft}{0pt}\setlength{\LTright}{0pt}
\setlength{\LTpre}{6pt}\setlength{\LTpost}{6pt}
\AtBeginEnvironment{xltabular}{\small}
% Bảng ngắn: không ngắt trang giữa bảng
\newenvironment{shorttable}{\par\nopagebreak\vspace{6pt}\nopagebreak\noindent\small\tabularx}{\endtabularx\par\vspace{6pt}}
% Vị trí hình: ưu tiên ngay tại chỗ, tránh khoảng trắng lớn
\renewcommand{\topfraction}{0.9}
\renewcommand{\bottomfraction}{0.8}
\renewcommand{\textfraction}{0.07}
\renewcommand{\floatpagefraction}{0.85}
\setcounter{topnumber}{3}\setcounter{bottomnumber}{2}\setcounter{totalnumber}{4}
\setlist{itemsep=1pt,topsep=3pt,parsep=0pt}

\definecolor{brand}{HTML}{1E3A8A}
\definecolor{muted}{HTML}{6B7280}

\titleformat{\section}{\needspace{8\baselineskip}\sffamily\Large\bfseries\color{brand}}{\thesection}{0.6em}{}
\titleformat{\subsection}{\needspace{6\baselineskip}\sffamily\large\bfseries}{\thesubsection}{0.6em}{}
\titleformat{\subsubsection}{\sffamily\normalsize\bfseries}{\thesubsubsection}{0.5em}{}
\titlespacing*{\section}{0pt}{12pt}{6pt}
\titlespacing*{\subsection}{0pt}{10pt}{4pt}

\pagestyle{fancy}
\fancyhf{}
\fancyfoot[C]{\small\sffamily\thepage}
\renewcommand{\headrulewidth}{0pt}

% \code{...}: định danh trong mã, giữ khoảng trắng, ngắt dòng tại / . _ - : , = |
\ExplSyntaxOn
\NewDocumentCommand \code { m }
  {
    \texttt
      {
        \str_map_inline:nn {#1}
          {
            ##1
            \str_case:nnF {##1}
              { {/}{\allowbreak} {.}{\allowbreak} {_}{\allowbreak}
                {-}{\allowbreak} {:}{\allowbreak} {,}{\allowbreak}
                {=}{\allowbreak} {|}{\allowbreak} }
              { }
          }
      }
  }
\ExplSyntaxOff

\newcommand{\ra}{\ensuremath{\rightarrow}}
\newcommand{\src}[1]{{\small\color{muted}Mã nguồn: \code{#1}}\par}

\newcolumntype{L}[1]{>{\raggedright\arraybackslash}p{#1}}
\newcolumntype{Y}{>{\raggedright\arraybackslash}X}

% Code block nhỏ
\DefineVerbatimEnvironment{codeblock}{Verbatim}{fontsize=\small,frame=leftline,
  framerule=1.2pt,rulecolor=\color{muted},xleftmargin=6pt,framesep=6pt}

% \figslot{tên-file}{chiều cao ô}{chú thích}
%   Có figures/<tên-file>.pdf|png|jpg thì chèn hình; chưa có thì để ô trắng.
%   Tham chiếu: \ref{fig:<tên-file>}
\usepackage{graphicx,float,caption}
\captionsetup{font=small,labelfont=bf,justification=centering,skip=6pt}
\graphicspath{{figures/}}
\newcommand{\figimg}[1]{\includegraphics[width=\linewidth,height=0.6\textheight,keepaspectratio]{#1}}
\newcommand{\figslot}[3]{%
  \begin{figure}[!htbp]\centering
    \IfFileExists{figures/#1.pdf}{\figimg{#1.pdf}}{%
    \IfFileExists{figures/#1.png}{\figimg{#1.png}}{%
    \IfFileExists{figures/#1.jpg}{\figimg{#1.jpg}}{%
      {\setlength{\fboxsep}{0pt}\color{muted!60}\fbox{\parbox[b][#2][b]{\dimexpr\linewidth-0.8pt}{%
        \hfill{\scriptsize\ttfamily\color{muted!80}figures/#1\ }\par\vspace{2pt}}}}}}}
    \caption{#3}\label{fig:#1}
  \end{figure}}


\begin{document}

{\sffamily
{\Large\bfseries\color{brand} Báo cáo kỹ thuật}\\[2pt]
{\large Hệ thống Marketplace Price Intelligence}\\[6pt]
\small
\begin{tabular}{@{}ll@{}}
Sinh viên: & \StudentName\ (MSSV \StudentId) \\
GVHD: & \Supervisor \\
\end{tabular}}

\tableofcontents

% =====================================================================
\section{Tổng quan hệ thống}
% =====================================================================
\subsection{Mục tiêu và phạm vi}
Hệ thống thu thập định kỳ trang danh mục của sàn TMĐT (Tiki là nguồn chính, Lazada dùng để kiểm tra khả năng thêm nguồn), lưu lại toàn bộ
dữ liệu gốc, rồi dựng kho dữ liệu phục vụ hai nhu cầu: phân tích lịch sử giá theo giờ/ngày (biến động giá, khuyến mãi đáng ngờ, độ phủ nguồn)
và phát hiện thay đổi giá gần thời gian thực. Shopee nằm ngoài phạm vi vì có anti-bot (ADR-0006).

\subsection{Kiến trúc tổng thể}
Kiến trúc Lambda, nạp theo kiểu ELT với ba tầng Bronze/Silver/Gold.
Crawler ghi dữ liệu vào Bronze trước, sau đó mới đẩy sang Kafka; từ Bronze tách ra hai nhánh xử lý:
\begin{itemize}
  \item \textbf{Batch} (mỗi giờ): Bronze \ra\ Silver (Delta) \ra\ star schema \code{dw} \ra\ 12 mart \ra\ kiểm tra chất lượng \ra\ publish vào PostgreSQL \ra\ promote.
  \item \textbf{Speed}: Kafka \code{records.v1} \ra\ chuẩn hoá \ra\ so sánh với trạng thái từng offer \ra\ sự kiện thay đổi ra Kafka, Elasticsearch, Redis.
\end{itemize}

\figslot{fig-1-kien-truc-tong-the}{80mm}{Kiến trúc tổng thể: nguồn, Bronze, nhánh batch, nhánh speed, lớp phục vụ}

\begin{xltabular}{\linewidth}{@{}L{32mm}L{50mm}Y@{}}
  \toprule
  \textbf{Thành phần} & \textbf{Công nghệ} & \textbf{Vai trò} \\
  \midrule
  Thu thập & Python, requests, Playwright & crawl trang danh mục, frontier trong PostgreSQL \\
  Data lake & MinIO (S3); S3 trên Databricks & Bronze, Silver, Gold \\
  Định dạng bảng & Delta Lake & Silver, Gold, time travel theo version \\
  Xử lý & Spark 4 (batch + Structured Streaming) & chuẩn hoá, mô hình, mart, speed \\
  Hàng đợi & Kafka (5 topic) & đường dữ liệu cho speed \\
  Điều phối & Airflow 3.3.2, DockerOperator & 4 DAG \\
  Phục vụ & PostgreSQL \code{cache.*}, Superset; Elasticsearch, Kibana, Redis & batch view; realtime view \\
  Giám sát & Alloy, Prometheus, Loki, Grafana & metrics, log, cảnh báo \\
  \bottomrule
\end{xltabular}

\subsection{Nguyên tắc thiết kế}
\begin{itemize}
  \item Lake là nơi lưu chính. PostgreSQL chỉ giữ cache của run đang phục vụ và trạng thái vận hành; Kafka, Elasticsearch, Redis đều dựng lại được từ Bronze.
  \item Mọi ID sinh tất định từ khoá tự nhiên (SHA-256, \code{xxhash64}), nên ghi lại hoặc replay không tạo bản trùng. Kafka dùng at-least-once.
  \item Mỗi run batch ghi ra thư mục riêng và chỉ được phục vụ sau khi qua cổng chất lượng và promote bằng compare-and-swap con trỏ \code{current.json}.
\end{itemize}

% =====================================================================
\section{Thu thập dữ liệu}
% =====================================================================
\src{src/ecomdata/ingestion/, infra/http.py, domain/sources.py}

\subsection{Nguồn}
\begin{xltabular}{\linewidth}{@{}L{30mm}YY@{}}
  \toprule
  & \textbf{Tiki} & \textbf{Lazada} \\
  \midrule
  Danh sách danh mục & sitemap \ra\ URL \code{<slug>/c<id>} & 20 danh mục khai báo tĩnh \\
  Cách tải & HTTP GET (\code{requests}) & Playwright Chromium, bắt XHR \code{ajax=true} \\
  Vị trí dữ liệu & JSON \code{__NEXT_DATA__} trong HTML & \code{mods.listItems} trong XHR \\
  Phân trang & chỉ trang 1 (40 sản phẩm) & tối đa 102 trang \\
  Delay tối thiểu & 2 s/host & 5 s/host \\
  Nhận biết bị chặn & không có \code{__NEXT_DATA__} & URL/HTML chứa \code{punish}, \code{captcha} \\
  \bottomrule
\end{xltabular}

Mọi request đi qua robots.txt (cache 1 giờ; không đọc được thì coi như cấm), throttle theo host (delay + jitter 0--1 s) và user-agent riêng của dự án.

\subsection{Frontier và lập lịch}
Hàng đợi việc là bảng \code{audit.crawl_frontier}; mỗi dòng là một lần crawl của một target (danh mục $\times$ trang) tại thời điểm \code{scheduled_for}.
Worker lấy tối đa 10 task mỗi lần:

\begin{codeblock}
WITH due AS (
  SELECT task_id FROM audit.crawl_frontier
  WHERE ((status IN ('READY','RETRY_WAIT') AND scheduled_for <= :now)
      OR (status = 'LEASED' AND lease_expires_at <= :now))
  ORDER BY priority DESC, scheduled_for ASC
  LIMIT :limit
  FOR UPDATE SKIP LOCKED)
UPDATE audit.crawl_frontier f
SET status='LEASED', lease_owner=:worker, lease_expires_at=:now+300s,
    attempts = f.attempts + 1
FROM due WHERE f.task_id = due.task_id
RETURNING ...;
\end{codeblock}

Khi task thành công, cùng transaction đó chèn task kế tiếp với \code{scheduled_for = completed_at + cadence}. Cadence theo tier:
\code{ACTIVE} 1 h, \code{NORMAL} 4 h, \code{COLD} 12 h, \code{SWEEP} 24 h. DAG \code{source_discovery} seed các target SWEEP mỗi ngày.

\figslot{fig-2-luong-crawl}{70mm}{Luồng crawl: frontier, worker, Bronze, Kafka, audit}

\subsection{Một lượt tải trang}
URL \ra\ robots \ra\ throttle \ra\ fetch \ra\ ghi raw \ra\ kiểm HTTP status \ra\ trích xuất \ra\ bỏ PII \ra\ ghi records \ra\ publish Kafka (\code{acks=all}) \ra\ ghi audit.
Raw được ghi trước bước kiểm status nên phản hồi 429/503 vẫn được lưu lại.

\subsection{Xử lý lỗi}
\begin{itemize}
  \item Retry: \code{RATE_LIMITED}, \code{TRANSIENT_NETWORK}, \code{SERVER_ERROR}, \code{BLOCKED}, \code{STORAGE_ERROR}, \code{PUBLISH_ERROR}. Không retry: \code{ROBOTS_DENIED}, \code{PARSE_ERROR}, \code{VALIDATION_ERROR}, \code{CLIENT_ERROR}.
  \item Backoff theo \code{Retry-After} nếu có, nếu không $30 \cdot 2^{n-1}$ s (trần 1800 s, jitter $\pm 20\%$), tối đa 5 lần.
  \item Circuit breaker theo nguồn: 5 lỗi retry liên tiếp thì mở 900 s; trong thời gian mở, task được trả lại hàng đợi mà không gọi site.
\end{itemize}

% =====================================================================
\section{Lưu trữ và định dạng dữ liệu}
% =====================================================================
\begin{xltabular}{\linewidth}{@{}L{25mm}L{39mm}L{25mm}L{25mm}Y@{}}
  \toprule
  \textbf{Tầng} & \textbf{Vị trí} & \textbf{Định dạng} & \textbf{Partition} & \textbf{Cách ghi, lưu giữ} \\
  \midrule
  \endfirsthead
  \toprule \textbf{Tầng} & \textbf{Vị trí} & \textbf{Định dạng} & \textbf{Partition} & \textbf{Cách ghi} \\ \midrule
  \endhead
  Bronze/raw & \code{ecommerce-bronze/marketplace/raw/} & \code{.bin.gz} + \code{metadata.json} & marketplace, ngày, giờ, crawl run & append-only, giữ 30 ngày \\
  Bronze/records & \code{.../marketplace/records/} & JSONL.gz, 1 file/trang & source, ngày, giờ & giữ vĩnh viễn \\
  Bronze/\allowbreak crawl\_audit & \code{.../marketplace/crawl_audit/} & JSONL.gz & ngày, giờ & bản sao audit crawl \\
  Kafka & 5 topic, 3 partition & JSON & key theo nguồn/offer & at-least-once, giữ 7 ngày \\
  Silver & \code{ecommerce-silver/marketplace/} & Delta & \code{source_code}, \code{observed_date} & ghi đè partition ngày trong cửa sổ \\
  Gold \code{dw} & \code{ecommerce-gold/marketplace/dw/} & Delta & fact theo \code{date_key} & \code{MERGE} theo khoá \\
  Gold runs & \code{.../runs/run_id=<run>/} & Delta & theo ngày & một thư mục mỗi run \\
  Manifest & \code{manifests/}, \code{current.json} & JSON & --- & compare-and-swap \\
  Cache & PostgreSQL \code{cache.marketplace_*} & bảng & --- & TRUNCATE + INSERT, 1 transaction \\
  Speed & Elasticsearch, Redis \code{rt:*} & document, hash & --- & upsert theo id; Redis TTL 24 h \\
  \bottomrule
\end{xltabular}

\figslot{fig-3-du-lieu-qua-cac-tang}{65mm}{Dữ liệu qua các tầng: raw \ra\ records \ra\ Silver \ra\ star \ra\ mart \ra\ cache}

Bronze/records giữ nguyên payload của sàn, chỉ bọc thêm thông tin lineage (rút gọn):
\begin{codeblock}
{"schema_version":"marketplace-record.v1", "record_id":"rec_9f2c...",
 "source_code":"tiki", "record_kind":"listing_row", "record_index":7,
 "raw_artifact_id":"raw_51ab...", "raw_uri":"s3a://ecommerce-bronze/.../part-0.bin.gz",
 "body_sha256":"e3b0...", "crawl_run_id":"crawl_run_...",
 "request_target":"dien-thoai-may-tinh-bang/c1789", "request_page":1,
 "fetched_at":"2026-11-14T03:12:45Z", "http_status":200,
 "payload":{"id":..., "price":..., "list_price":..., "discount_rate":..., ...}}
\end{codeblock}
Tầng này không đổi tên, ép kiểu hay lọc trường (ELT, ADR-0001), nên có thể chuẩn hoá lại từ đầu khi đổi logic.

% =====================================================================
\section{Chuẩn hoá: Bronze \ra\ Silver}
% =====================================================================
\src{src/ecomdata/processing/normalize/, processing/batch/stages/silver.py}

Mỗi run dựng lại Silver cho cửa sổ gồm ngày của \code{as_of} và 1 ngày trước đó (để nhận dữ liệu đến muộn), thay toàn bộ các partition ngày đó.
Các bước:
\begin{enumerate}
  \item Bỏ record trùng theo \code{record_id}.
  \item Kiểm envelope (\code{schema_version}, nguồn đã đăng ký); sai thì vào quarantine.
  \item Normalizer của từng nguồn chuyển \code{listing_row} sang dạng canonical và dạng wide theo nguồn.
  \item Tách quarantine theo lý do: thiếu listing id, trùng trên trang, thiếu tiêu đề, giá thiếu/sai, thiếu \code{fetched_at}, thiếu lineage.
  \item Ghi 4 nhóm bảng Delta: \code{offer_observations_v2} (53 cột), \code{tiki_listing_wide}/\code{lazada_listing_wide}, \code{page_context}, \code{quarantine}.
\end{enumerate}

\begin{xltabular}{\linewidth}{@{}L{38mm}L{44mm}Y@{}}
  \toprule
  \textbf{Cột canonical} & \textbf{Trường nguồn (Tiki)} & \textbf{Quy tắc} \\
  \midrule
  \code{current_price} & \code{price} & \code{decimal(38,6)}, $\ge 0$ \\
  \code{list_price} & \code{list_price} / \code{original_price} & giữ nếu $> 0$ \\
  \code{discount_percent} & \code{discount_rate} & ngoài $[0,100]$ \ra\ null \\
  \code{sold_count} & \code{quantity_sold.value} & số nguyên $\ge 0$ \\
  \code{availability} & \code{availability}, \code{shippable} & \code{IN_STOCK} / \code{OUT_OF_STOCK} / \code{UNKNOWN} \\
  \code{is_advertisement} & \code{advertisement} & giữ dòng, gắn cờ \\
  \code{ranking_position} & thứ tự dòng organic & $(\text{page}-1)\cdot 40 + \text{rank}$ \\
  \code{offer_id} & sha256(code, listing id) & ổn định qua các lần crawl \\
  \code{observation_id} & sha256(code, listing id, observed\_at, body\_sha256) & duy nhất mỗi lần thấy \\
  \code{observed_date} & \code{to_date(observed_at)} & theo session timezone UTC \\
  \bottomrule
\end{xltabular}

% =====================================================================
\section{Mô hình dữ liệu Gold}
% =====================================================================
\src{src/ecomdata/processing/batch/ (dimensions.py, facts.py, marts.py, discount.py)}

\subsection{Star schema \texttt{dw}}
\begin{itemize}
  \item 7 dimension: \code{dim_date}, \code{dim_marketplace}, \code{dim_seller}, \code{dim_brand}, \code{dim_category}, \code{dim_product_listing}, \code{dim_listing_flags};
        1 bridge \code{bridge_category_ancestry}.
  \item 6 fact: \code{fact_offer_observation}, \code{fact_price_snapshot_daily}, \code{fact_price_change}, \code{fact_counter_daily}, \code{fact_category_page_daily}, \code{fact_crawl_attempt}.
  \item Surrogate key \code{xxhash64} của khoá tự nhiên; giá trị thiếu dùng khoá $-1$.
  \item \code{dim_product_listing} là SCD2 theo tiêu đề, brand, category, seller.
  \item Ghi bằng \code{MERGE} (update / insert / delete khi không còn ở nguồn) từ toàn bộ Silver.
\end{itemize}

\figslot{fig-4-star-schema}{90mm}{Star schema \texttt{dw}: fact, dimension và bridge}

\subsection{Mart}
Stage \code{marts} đọc star tại đúng version Delta vừa ghi, lọc \code{observed_at <= as_of}, bỏ dòng quảng cáo và dựng 12 mart:
\code{offer_current}, \code{seller_current}, \code{offer_price_history_daily}, \code{offer_change_daily}, \code{offer_freshness}, \code{category_price_daily},
\code{source_coverage_daily}, \code{crawl_reliability_daily}, \code{counter_delta_daily}, \code{price_anomaly_daily}, \code{discount_credibility_daily}, \code{discount_daily}.
Manifest của run ghi version Delta của từng bảng star để tái lập.

\subsection{Logic phân tích}
\begin{itemize}
  \item Giá bất thường: robust z-score so với 14 ngày trước của chính offer (cần $\ge 7$ ngày), ngưỡng 3,5; MAD $= 0$ thì dùng IQR.
  \item Độ tin cậy khuyến mãi: giá tham chiếu là giá thấp nhất 30 ngày trước; 5 luật có trọng số (giá cao hơn giá niêm yết, giảm $>50\%$, tăng giá niêm yết,
        mức giảm công bố vượt lịch sử, giá không thấp hơn tham chiếu); tổng điểm $\ge 50$ là SUSPECT, $\ge 20$ là QUESTIONABLE. Isolation Forest chỉ dùng để xếp hạng.
  \item Counter (đã bán, đánh giá): mỗi bước chuyển gắn nhãn \code{VALID}, \code{GAP_TOO_LONG}, \code{COUNTER_DECREASED}; không suy ra số đơn hàng.
\end{itemize}

% =====================================================================
\section{Speed layer}
% =====================================================================
\src{src/ecomdata/processing/speed/, domain/change_rules.py}

Một ứng dụng Spark Structured Streaming chạy hai query, trigger 30 s:
\begin{enumerate}
  \item \textbf{Normalizer}: \code{records.v1} \ra\ dùng lại normalizer của batch \ra\ bỏ quảng cáo \ra\ \code{observations.v1}; dòng lỗi \ra\ \code{records.v1.dlq}.
  \item \textbf{Stateful}: kiểm hợp đồng (lỗi \ra\ \code{observations.v1.dlq}), \code{groupBy(offer_id)} + \code{applyInPandasWithState}, state là quan sát mới nhất của offer.
        Bản trùng và bản đến muộn bị bỏ; bản mới được so với state để sinh sự kiện.
\end{enumerate}

Sự kiện: \code{NEW_OFFER}, \code{PRICE_CHANGED}, \code{LARGE_PRICE_DROP} ($\ge 100.000$đ hoặc $\ge 20\%$), \code{RATING_CHANGED}, \code{COUNTER_CHANGED},
\code{AVAILABILITY_CHANGED}, \code{OFFER_STALE} (6 giờ không thấy). Ghi lần lượt ra Kafka \code{changes.v1}, Elasticsearch, Redis, audit;
micro-batch đã ghi thành công thì bỏ qua khi replay.

\figslot{fig-5-speed}{60mm}{Speed layer: hai query, topic Kafka, DLQ và các sink}

% =====================================================================
\section{Chất lượng dữ liệu và phát hành}
% =====================================================================
\src{src/ecomdata/domain/quality_rules.py, processing/batch/stages/ (quality.py, publish.py, promote.py)}

\begin{xltabular}{\linewidth}{@{}L{30mm}Y@{}}
  \toprule
  \textbf{Nhóm luật} & \textbf{Nội dung} \\
  \midrule
  Lineage & checksum, URI của raw tồn tại \\
  Duy nhất & \code{observation_id}, khoá listing, một dòng/offer ở mart hiện hành \\
  Miền giá trị & giá $\ge 0$, tiền tệ hợp lệ, \code{observed_at} không ở tương lai \\
  Đối soát & Silver + quarantine = \code{parsed_count} của audit; số dòng, tổng giá khớp giữa fact và mart \\
  Star & khoá ngoại hợp lệ, surrogate key duy nhất, SCD2 liên tục và có đúng 1 version hiện hành \\
  \bottomrule
\end{xltabular}

Tổng 22 luật, 16 bắt buộc. Một luật bắt buộc không đạt thì run thành \code{QUALITY_FAILED} và star được \code{RESTORE} về version của run đang phục vụ.
Run đạt thì:
\begin{itemize}
  \item \textbf{Publish}: ghi các mart vào \code{staging.*}, sau đó trong một transaction (có advisory lock) TRUNCATE và INSERT 12 bảng \code{cache.*}.
  \item \textbf{Promote}: compare-and-swap \code{current.json}; nếu con trỏ đã bị run khác đổi thì run chuyển sang \code{HELD}.
\end{itemize}
Đã kiểm tra trường hợp lỗi: run 07/11 11:00 bị \code{QUALITY_FAILED}, Superset vẫn giữ dữ liệu của run trước.

% =====================================================================
\section{Điều phối}
% =====================================================================
\src{src/ecomdata/orchestration/dags/, processing/batch/runner.py}

\begin{xltabular}{\linewidth}{@{}L{48mm}L{22mm}Y@{}}
  \toprule
  \textbf{DAG} & \textbf{Lịch} & \textbf{Nội dung} \\
  \midrule
  \code{marketplace_batch} & hằng giờ & sensor 10 phút, rồi 7 stage tuần tự \\
  \code{marketplace_reprocess} & thủ công & chạy lại một khoảng hoặc rebuild toàn bộ \\
  \code{source_discovery} & hằng ngày & tìm danh mục, seed frontier SWEEP \\
  \code{marketplace_maintenance} & hằng tuần & dọn raw quá hạn, bảo trì \\
  \bottomrule
\end{xltabular}

Mỗi stage của \code{marketplace_batch} là một task \code{DockerOperator} chạy container Spark riêng:
\code{silver} \ra\ \code{dimensions} \ra\ \code{facts} \ra\ \code{marts} \ra\ \code{quality} \ra\ \code{publish} \ra\ \code{promote}.
\code{max_active_runs=1}, retry 2 lần (trừ \code{quality}), timeout 45 phút/stage. Trạng thái run lưu ở \code{audit.marketplace_batch_run};
stage đã hoàn thành được bỏ qua khi chạy lại. Advisory lock chặn hai batch chạy song song.

\figslot{fig-6-dag-batch}{45mm}{DAG \texttt{marketplace\_batch} và các trạng thái của run}

% =====================================================================
\section{Giám sát}
% =====================================================================
\src{observability/ (alloy/, grafana/dashboards/, grafana/provisioning/alerting/rules.yaml)}

\begin{itemize}
  \item Grafana Alloy thu log container và log task Airflow vào Loki, metrics vào Prometheus.
  \item Metrics hạ tầng: cAdvisor, Kafka consumer lag, Redis, Elasticsearch, StatsD của Airflow, probe HTTP.
  \item Metrics nghiệp vụ: Alloy chạy truy vấn SQL trên các bảng \code{audit.*} mỗi phút, sinh các nhóm \code{mp_crawl_*}, \code{mp_batch_*}, \code{mp_quality_*},
        \code{mp_speed_*}, \code{mp_serving_cache}, \code{mp_storage}.
\end{itemize}

\begin{codeblock}
-- mp_crawl_5m: số lượt crawl 5 phút gần nhất theo nguồn, trạng thái, loại lỗi
SELECT f.marketplace_code AS marketplace, a.status,
       COALESCE(a.error_kind,'NONE') AS error_kind,
       count(*) AS attempts, sum(a.raw_bytes) AS raw_bytes
FROM audit.crawl_request_attempt a JOIN audit.crawl_frontier f USING (task_id)
WHERE a.completed_at > now() - interval '5 minutes'
GROUP BY 1,2,3;
\end{codeblock}

\figslot{fig-7-giam-sat}{55mm}{Luồng tín hiệu giám sát: nguồn metrics/log, Alloy, Prometheus, Loki, Grafana}

Grafana có 5 dashboard (overview, pipeline, resources, orchestration, logs) và 19 luật cảnh báo:
\begin{xltabular}{\linewidth}{@{}L{22mm}Y@{}}
  \toprule
  \textbf{Nhóm} & \textbf{Cảnh báo} \\
  \midrule
  Hạ tầng & endpoint không phản hồi; RAM $>90\%$; đĩa $>85\%$; OOM kill; container khởi động lại liên tục; Alloy mất metrics \\
  Pipeline & crawl ngừng; tỉ lệ bị chặn $>20\%$; circuit mở; speed không có micro-batch thành công; lag $>5000$ offset; không có run promote trong 2,5 h;
             batch lỗi hoặc HELD; luật quality bắt buộc FAIL; DLQ $>100$/30 phút; consumer lag $>10.000$; nhiều log lỗi \\
  Điều phối & DAG run gần nhất lỗi; scheduler mất heartbeat $>120$ s \\
  \bottomrule
\end{xltabular}

% =====================================================================
\section{Kết quả vận hành}
% =====================================================================
Lượt chạy toàn luồng 13--14/11/2026 trên một laptop (\code{docs/evidence/fullrun-20261114/}):

\begin{xltabular}{\linewidth}{@{}L{55mm}Y@{}}
  \toprule
  \textbf{Chỉ số} & \textbf{Giá trị} \\
  \midrule
  Thời gian chạy & 9 giờ 16 phút (dừng khi RAM host còn 2,6 GB) \\
  Run batch & 10/10 PROMOTED, quality PASS \\
  Số dòng Gold mỗi run & 1.432.017 \ra\ 1.966.878 \\
  Thời gian stage & \code{silver} 38--189 s; \code{dimensions} 217--259 s; \code{marts} 88--105 s \\
  Crawl & khoảng 1.270 trang/giờ trong 4 giờ đầu, sau đó 25 trang/giờ \\
  Speed & 0 lỗi, DLQ 0; độ trễ p95 59--62 s sau khi hết tồn đọng \\
  Khuyến mãi đáng ngờ & 34,9\% dòng có giảm giá bị gắn cờ (ngày 13/11) \\
  \bottomrule
\end{xltabular}

\figslot{chart-1-crawl-theo-gio}{55mm}{Số trang crawl thành công theo giờ, 13--14/11}
\figslot{chart-2-thoi-gian-stage}{55mm}{Thời gian từng stage batch qua các run}

% =====================================================================
\section{Hạn chế hiện tại}
% =====================================================================
\begin{xltabular}{\linewidth}{@{}L{22mm}YY@{}}
  \toprule
  \textbf{Phần} & \textbf{Hạn chế} & \textbf{Hướng xử lý} \\
  \midrule
  \endfirsthead
  \toprule \textbf{Phần} & \textbf{Hạn chế} & \textbf{Hướng xử lý} \\ \midrule
  \endhead
  Thu thập & Tiki chỉ lấy trang 1 nên mẫu là nhóm sản phẩm sàn đẩy lên đầu & ghi rõ phạm vi kết luận \\
  Thu thập & Lịch SWEEP dồn vào khoảng 4 giờ rồi nghỉ khoảng 20 giờ & rải \code{scheduled_for} trong ngày \\
  Thu thập & \code{attempts} tăng cả khi task bị trả về do circuit/lease hết hạn, có thể vượt \code{max_attempts} và làm lỗi cả lô lease & chỉ lease khi \code{attempts < max_attempts} \\
  Lưu trữ & Records 1 file/trang, khoảng 10 KB/file & compaction theo giờ \\
  Lưu trữ & Chưa có \code{OPTIMIZE}/\code{VACUUM}, thư mục run Gold cũ không bị xoá & thêm vào DAG maintenance \\
  Chuẩn hoá & Ngày nghiệp vụ tính theo UTC, lệch 7 giờ so với ngày ở Việt Nam & tính cột ngày theo \code{Asia/Ho_Chi_Minh} \\
  Chuẩn hoá & Quarantine chỉ lưu lý do đầu tiên & lưu mảng lý do \\
  Gold & Dimension, fact, mart dựng lại từ toàn bộ lịch sử mỗi giờ, thời gian tăng theo số ngày dữ liệu & fact incremental theo partition ngày \\
  Gold & Luật khuyến mãi gắn cờ 34,9\% khi mới có khoảng 4 ngày lịch sử & chỉ bật luật lịch sử khi đủ số ngày; gán nhãn tay để đánh giá \\
  Speed & State của offer không bị xoá khi timeout, tăng theo số offer từng thấy & \code{state.remove()} sau thời gian dài không thấy \\
  Speed & Sink lỗi tạm thời làm query dừng, phụ thuộc Docker khởi động lại & retry có backoff trong sink \\
  Phục vụ & Batch view và realtime view phục vụ riêng, chưa có view hợp nhất & --- \\
  Giám sát & Metrics nghiệp vụ phụ thuộc PostgreSQL vận hành; chưa có cảnh báo độ tươi theo nguồn & thêm cảnh báo theo nguồn \\
  \bottomrule
\end{xltabular}

\end{document}
